\subsection{معرفی و انگیزه}
این ماژول در چارچوب پژوهش حاضر پیاده‌سازی و ارزیابی شده و فاقد مقاله مستقل جداگانه در ادبیات تحقیق است.

مکانیزم توجه \lr{Kernelized Decay Attention (KDA)} (یا توجه خطی بازگشتی با زوال آموخته‌شده) با هدف ایجاد پیوند نظری میان مدل‌های توجه خطی با پیچیدگی $\mathcal{O}(N)$ و مکانیزم‌های توجه بازگشتی مقید به حافظه با سوگیری‌های القایی قوی بر مبنای فاصله توکن‌ها توسعه یافته است. در الگوهای سنتی ترنسفورمر، توجه کامل دوقطبی فارغ از فاصله زمانی میان نشانه‌ها، ظرفیت محاسباتی یکسانی را اختصاص می‌دهد که علاوه بر ایجاد سربار درجه دو، در بسیاری از کاربردها به پراکندگی توجه بر روی توکن‌های نویزی و نامرتبط منجر می‌گردد.

ایده کلیدی در \lr{KDA} معرفی یک ضریب زوال نمایی یادگیرنده (\lr{Learned Exponential Decay Gate}) در فرمولاسیون توجه هسته‌محور است. این زوال پیوسته به شبکه امکان می‌دهد تا به طور خودکار افق زمانی مؤثر (\lr{Effective Context Window}) را برای هر سر توجه تنظیم نماید. علاوه بر این، با استفاده از بهنجارسازی کروی، ردیابی مقادیر مخرج به صورت برخط و نگاشت فاکتورگیری‌شده \lr{LoRA} برای انتقال وزن‌ها، پایداری عددی گرادیان‌ها در توالی‌های قابل توجه طولانی تضمین می‌شود.
